% META: {"id":"1SPE-EXPO-CO-028","chapitre":"1SPE-EXPONENTIELLE","type_objet":"corrige","exercice_id":"1SPE-EXPO-EX-028","status":"approved","extension_codes":["X1"],"programme_alignment":"OPTIONAL_EXTENSION","extension_label":"Approfondissement — Vers la Terminale","origin":"generated","status_history":["generated","audited","accepted","approved"],"release_acceptance":"RELEASE_OWNER_BATCH_ACCEPTANCE","acceptance_closure_digest":"sha256:9b3ccf9a81c5520fbb7e03b7b2d3e7bf2057a02834b6d908bf3be5f49f3b553f","acceptance_receipt_digest":"sha256:4fad86f0282e0fa4e0419cca1ad6379ae7af5a280c04a4a90880f90a1c044242"}
\begin{center}\textbf{Approfondissement — Vers la Terminale}\end{center}
% BEGIN-VERIFY
% from sympy import *
% x = symbols('x')
% g = exp(x) - x - 2
% assert g.subs(x, 0) == -1
% assert g.subs(x, 2) == exp(2) - 4
% assert g.subs(x, 0) < 0
% assert g.subs(x, 2) > 0
% END-VERIFY
\begin{corrige}{1SPE-EXPO-EX-028}
\textbf{1.} On calcule :
\[g(0) = \mathrm{e}^0 - 0 - 2 = 1 - 2 = -1 < 0.\]
\[g(2) = \mathrm{e}^2 - 2 - 2 = \mathrm{e}^2 - 4 \approx 7{,}389 - 4 = 3{,}389 > 0.\]

\textbf{2.} La fonction $g$ est la somme de la fonction exponentielle (continue sur $\mathbb{R}$) et de la fonction affine $x \mapsto -x - 2$ (continue sur $\mathbb{R}$). Donc $g$ est \textbf{continue} sur $\mathbb{R}$, et en particulier sur $[0\,;\,2]$.

\textbf{3.} On a :
\begin{itemize}
  \item $g$ est continue sur $[0\,;\,2]$ ;
  \item $g(0) = -1 < 0$ et $g(2) = \mathrm{e}^2 - 4 > 0$.
\end{itemize}
D'après le \textbf{théorème des valeurs intermédiaires}, comme $0$ est compris entre $g(0)$ et $g(2)$, il existe au moins un réel $c \in [0\,;\,2]$ tel que $g(c) = 0$, c'est-à-dire $\mathrm{e}^c = c + 2$.

L'équation $\mathrm{e}^x = x + 2$ admet donc au moins une solution dans $[0\,;\,2]$.
\end{corrige}
